\section{The lattice source and rational normalization}\label{sec:source}

We first fix the finite input and its error budget. The output of this
section is a signed rational Heisenberg Hamiltonian with polynomial
geometric extent and coefficients, together with a positive
inverse-polynomial decision gap.

Our starting problem is the field-free square-lattice Heisenberg energy
problem, which is QMA-hard by Theorem~48 and Section~10.1 of Cubitt,
Montanaro and Piddock \cite{CMP2019}, using the polynomial-weight
spatially sparse construction of their Lemma~47. The Hamiltonian on a
finite, nonwrapping square-lattice subgraph is
\[
             H_{\rm in}=\sum_e J_e h^s_{LR},\qquad
             h^s_{LR}=\boldsymbol\sigma_L\cdot\boldsymbol\sigma_R
\]
on spins \(1/2\), where the components are Pauli matrices and there are
no one-site terms. YES instances have ground energy at most the lower
threshold, and NO instances have ground energy at least the upper
threshold. The cited construction supplies polynomial bounds on lattice
size and coefficient magnitudes and an inverse-polynomial promise gap.
The following deterministic normalizations give signed rational
coefficients, rational thresholds \(a_s,b_s\), and polynomial geometric
extent.

For completeness, let \(\delta>0\) be a known rational
inverse-polynomial lower bound for the incoming promise gap, and let
\(q\) be the number of weighted edges. Set
\(\varepsilon_0=\delta/100\). Rounding each weight within
\(\varepsilon_0/\max(1,q)\) changes the Hamiltonian norm by at most
\(3\varepsilon_0\). Write the incoming scalar as \(c_0I\), with
\(c_0=0\) if absent, and the incoming thresholds as \(a_0,b_0\).
Approximate these three numbers within \(\varepsilon_0\), obtaining
\(\widetilde c_0,\widetilde a_0,\widetilde b_0\). The rational
Hamiltonian without the scalar then has valid thresholds
\[
 a_s=\widetilde a_0-\widetilde c_0+5\varepsilon_0,\qquad
 b_s=\widetilde b_0-\widetilde c_0-5\varepsilon_0.
\]
The offsets allow \(3\varepsilon_0\) for the Hamiltonian,
\(\varepsilon_0\) for the scalar, and \(\varepsilon_0\) for the
threshold. Thus both promise implications survive, and
\(b_s-a_s\ge\delta-12\varepsilon_0>0\). Polynomial magnitudes and
inverse-polynomial accuracy require only polynomial rational bit lengths.
A connected square-grid component with \(V\) vertices has coordinate
diameter at most \(V-1\), by following lattice paths. Translating
components into disjoint strips therefore gives polynomial extent
without changing their weighted graphs.

Put \(n=\max(2,\text{source input size})\), \(\gamma=\min(1,b_s-a_s)\). All polynomial choices below can use fixed size, magnitude and inverse-gap bounds for this source. We can translate coordinates to keep them polynomially bounded, omit zero weights, and assume a nonempty vertex set (add a free spin if needed).

The quantity $\gamma$ is the error budget used in the finite reductions
below. It concerns the normalized spin instance; the final electronic
thresholds will be obtained from an explicitly computed affine energy
map in Section~\ref{subsec:physical-thresholds}.
